% RESEARCH_DRAFT_SHA256: 1113a227d247812c4a5d9cf5708481a67b901ebd04a9dbd5221df22902ec3942
% Layout and source-reference wrappers expanded for the web renderer.
\documentclass[11pt]{article}
\usepackage[margin=25mm]{geometry}
\usepackage{fontspec,unicode-math}
\setmainfont{Latin Modern Roman}
\setmathfont{Latin Modern Math}
\usepackage{amsmath,amssymb,mathtools,xurl,hyperref}
\hypersetup{colorlinks=true,urlcolor=blue,linkcolor=blue}
\setcounter{secnumdepth}{0}
\setlength{\parindent}{0pt}
\setlength{\parskip}{5pt}
\setlength{\emergencystretch}{3em}
\newcommand{\sref}[2]{\hyperlink{source:#1:#2}{[S#2]}}
\newcommand{\cataloguescope}[1]{\paragraph{Source record.}#1}
\newcommand{\floor}[1]{\left\lfloor#1\right\rfloor}
\begin{document}
\section*{Research-attempt provenance}
This standalone research attempt retains the complete original definition body below,
including its historical source and preparation notes. Those notes describe the
earlier catalogue preparation; the new attempt and its verification are identified
separately after that body. No original ID, catalogue status, or source statement is changed.

% UnsolvedMath ID 30006388; source catalogue code OWR-14299517-007
\section{Iterative Quantum LDPC Constructions from Weaker Seeds}\label{30006388}
{\small UnsolvedMath ID 30006388 · Computer Science}
\subsection{Definitions and mathematical statement}

\paragraph{Shared notation.}Standard number systems are $\mathbb N,\mathbb Z,\mathbb Q,\mathbb R,\mathbb C$. Specialized notation and index ranges follow the cited source and the definitions below.


A code $C\subseteq\Sigma^N$ is a set of permitted length-$N$ words. Its Hamming minimum distance is the minimum number of unequal coordinates of two distinct words. For a linear $[N,K,D]_q$ code, rate is $K/N$ and relative distance is $D/N$; decoding recovers a codeword from corrupted observations.

A quantum LDPC code has checks of bounded weight with each qubit participating in only boundedly many checks. Asymptotic goodness means both the logical-qubit rate and relative distance stay bounded away from zero as block length grows.
\paragraph{Record scope and qualifications.}
The catalogue title is retained as metadata. Named certificates, contracts, bounds or frameworks in that title are not asserted as results of the original source extract.
The permitted seed objects, iteration, and meaning of weaker are not specified. This is not changed to the already broader existence question for good quantum LDPC codes.

Source-record limitation: The permitted seed object and meaning of weaker are unspecified; the iterative construction target is absent.
\paragraph{Statement recorded in the supplied source.}
Is there an iterative construction of asymptotically good quantum low-density parity-check codes that starts from a weaker seed object? \href{https://doi.org/10.4171/owr/2025/41}{[S2]}
\paragraph{Source status.}
The supplied catalogue labels this source record open; that metadata was checked against the supplied snapshot on 2026-10-11. The label is preserved for provenance and is not an independent verification of present open status or mathematical correctness. The catalogue literature assessment is dated 2026-08-22.
\subsection{Short English statement}
Iteratively build good quantum sparse-check codes from a weaker starting object.
\subsection{Sources}
\begin{itemize}
\item[S1] Ulam AI and the UnsolvedMath Contributors, UnsolvedMath: A Curated Collection of Open Mathematics Problems, record 30006388 (OWR-14299517-007). Catalogue curation and metadata: CC BY 4.0. The source extract is retained; definitions, English summaries and qualifications are editorial additions.
\url{https://huggingface.co/datasets/ulamai/UnsolvedMath}.
\item[S2] New Mathematical Directions in Coding Theory (2026), Oberwolfach Reports, DOI 10.4171/owr/2025/41, pp. 2187--2242
\url{https://doi.org/10.4171/owr/2025/41}.
\item[S3] Quantum LDPC Codes of Almost Linear Distance
\url{https://doi.org/10.4230/LIPIcs.CCC.2025.25}.
\item[S4] Oberwolfach Report 41 2025
\url{https://ems.press/content/serial-article-files/52245}.
\item Source-specific Oberwolfach Reports licences, including CC BY-SA 4.0 where applicable, are retained; see THIRD\_PARTY\_NOTICES.md.
\end{itemize}
\paragraph{Preparation and review.}
The original source and UnsolvedMath contributors retain attribution for the source extract. Codex (GPT-6; exact serving revision not exposed) assisted with definitions, English summaries, scope notes and formatting. No new human mathematical authorship or independently verified solution is asserted. Permanent source ID 30006388 and code OWR-14299517-007 checked against the complete supplied registry and current repository. Distinct upstream records, including repeated questions and broad or damaged extracts, are retained for complete coverage of the source CS category. A repeated record is not a new independent conjecture. Complete mathematical-equivalence and current-literature review has not been performed. Preparation and final review follow contribution prompts C and D, with source gaps disclosed. Structural and rendering checks do not establish a new theorem or open status.
\label{end:30006388}

\section{Literature and context check}
The primary report, printed pp. 2223--2224, already gives an iterative almost-good subsystem-code construction from a strong constant-sized quantum locally testable seed and asks for weaker starting objects. Its theorem and homological-product distance estimates are not re-proved here. The following concrete attempt starts instead from an arbitrary binary classical check matrix and shows exactly why a simple quantum conversion fails.
\section{Research attempt: A commuting-check construction with unavoidable weight-two logical operators}

Let $H$ be an $s\times n$ binary matrix of rank $r<n$. A CSS code uses binary check matrices $H_X,H_Z$ with $H_XH_Z^{\mathsf T}=0$: an $X$-type check on support $u$ and a $Z$-type check on support $v$ commute exactly when $u\cdot v=0$ over $\mathbb F_2$, because each overlapping position contributes a minus sign. Taking $H_X=H_Z=H$ generally fails this requirement, even for the one-row matrix $(1,1,1)$, whose self-dot-product is one.

A simple repair is to double the coordinates and set
\[
H_X=H_Z=[\,H\ H\,].
\]
Now the product is $HH^{\mathsf T}+HH^{\mathsf T}=0$. The two types of checks have $r$ independent generators each: products of $X$-only and $Z$-only generators can equal identity only when both support vectors vanish. Hence the code on $2n$ qubits has $2n-2r>0$ logical qubits. Row weights double and column weights are unchanged, so bounded classical check degrees give bounded quantum check degrees.

However this repair always creates a nontrivial logical operator of weight two. Since $\operatorname{row}(H)$ is a proper subspace of $\mathbb F_2^n$, some coordinate vector $e_i$ is not in it; otherwise all basis vectors and hence the entire space would be in the row space. Let $z=(e_i,e_i)$. Then $[H\ H]z=0$, so the $Z$-operator with support $z$ commutes with all $X$ checks and automatically with all $Z$ checks. It is not a stabilizer: a pure $Z$ stabilizer has support $(u,u)$ with $u\in\operatorname{row}(H)$, whereas $e_i$ is outside that row space. Thus the quantum distance is at most two, independently of $n$ and independently of the distance of the original classical code. This rules out this specific conversion as a route to asymptotically good quantum codes.

\paragraph{Why simply repeating a quantum seed also fails.}
For a tensor product of $t$ disjoint copies of any fixed stabilizer code encoding at least one qubit, the total physical and logical dimensions grow by a factor $t$, and check weights stay fixed. A minimum-weight nontrivial logical Pauli on one block, tensored with identities elsewhere, remains nontrivial and has the seed's constant weight. In fact every Pauli commuting with all block checks restricts to a commuting Pauli on each block; if each restriction is a stabilizer, the entire operator is a stabilizer. Hence the product distance is exactly the seed distance. Sparse local checks and positive rate alone do not prevent logical errors confined to one block. An iterative weak-seed construction needs a mechanism that excludes such operators while maintaining sparsity; neither doubling checks nor disjoint products supplies it.
\section{Verification and adversarial audit}
The CSS commutation relation, check ranks, positive logical dimension, and non-stabilizer condition are all verified directly. The condition $r<n$ is necessary: at full rank the doubled code has no logical qubits, so a distance obstruction for encoded information would be meaningless. The code-distance argument is an upper bound and does not claim that every doubled construction has exact distance two.
\paragraph{Audit corrections.} Checked that the displayed weight-two operator is outside the stabilizer group, rather than merely commuting with checks; excluded the zero-logical-dimension boundary.
\paragraph{External computation.} None; symbolic arguments and boundary checks only.
\section{Scope and remaining gaps}
No homological-product subsystem-distance estimate or weaker seed sufficient for the original iterative construction is established. The proof identifies concrete failed alternatives without claiming impossibility of the requested construction.
\paragraph{Integration review.} Codex (GPT-6; exact serving revision not exposed) reviewed the written arguments and expanded layout and source-reference wrappers. Repository metadata and the production web renderer were checked using Python 3.13 and JavaScript with MathJax 3.2.2. This model-assisted review is not independent human proof verification; no PDF compilation or formal proof checking is asserted.

\section{Completion Estimate}
\noindent\textbf{COMPLETION: 0\%}
The calculations below the preserved source statement investigate a concrete route,
but establish no substantiated new progress toward the entire original target.
Known special cases and failed sufficient conditions are not credited as a solution.
\section{References}
The original source references and catalogue attribution are preserved in the
Sources section above. The following URLs were actually consulted for this attempt;
retention of an additional original reference does not claim fresh consultation of it.
\begin{itemize}
\item \url{https://ems.press/content/serial-article-files/52245} (consulted 2026-10-11). 
\end{itemize}
\section{Model and contributor disclosure}
Codex GPT-6 completed written-proof review of this attempt at its declared scope; this is model-assisted review, not independent human proof checking.
OpenAI gpt-6-astra, requested reasoning configuration ultra, generated this research
attempt on 2026-10-11. The exact serving revision is not exposed. Codex GPT-6
performed orchestration and repository-format review; its exact serving revision
is also not exposed. The mathematical route, boundary checks and final adversarial
self-audit followed contribution prompts A then B in
\texttt{docs/contribution\_prompts.md}, after reading \texttt{CONTRIBUTING.md}.
Python 3.14.2 assembled the standalone documents; any mathematical computation is
specified in Verification. No theorem prover was used. No human mathematical
authorship, substantial human proof correction, or independent human proof checking
is asserted. Model-assisted review, source retrieval, and structural checks do not
establish correctness or novelty. Full private reasoning transcripts are unavailable
as reproducibility artifacts; the complete arguments, inputs and outcomes actually
used are stated here or in the companion batch audit.
\section{Final status}
Unresolved as a new research contribution. This is a concrete route investigation,
not a claim of new complete resolution or a current-status certification. A known
answer reproduced for context does not become a new contribution. The catalogue's
source-status label is preserved; this attempt status does not assert that a
source-labelled solved question is open.
\end{document}
